% 第 14 章。数字：checks/check_ch14.py；引文：data/snapshots/africa_money_facts.csv（逐字核对见
% data/fetch/verify_web_quotes.py；达尔泽尔、奥维多与伊本·白图泰、伯顿的几页按扫描页逐字读，见快照元数据）；
% 图：figsrc/fig_ch14.py（inca_store 是教学模型）。
\chapter{非洲与美洲：黄金、贝币与没有货币的帝国}\label{ch:africa}

\begin{chapterguide}
\guidequestion{货币一定要由国家发行吗？一个庞大的帝国能不能完全不用货币？}
\begin{guidegoals}
\item 描述跨撒哈拉黄金贸易与马里帝国（曼萨·穆萨 1324 年朝觐的记载）；
\item 解释贝币作为跨区域货币的运作，以及它与奴隶贸易的关联和近代的贬值；
\item 分析印加帝国以劳役税（米塔制）与仓储代替货币的经济组织；
\item 比较阿兹特克的可可豆货币，归纳货币出现与否的条件。
\end{guidegoals}
\guideprereq{\chref{ch:mesopotamia}\enspace\nameref{ch:mesopotamia}；\chref{ch:silver}\enspace\nameref{ch:silver}}
\end{chapterguide}

伊斯兰历 724 年（1324），马里的国王曼萨·穆萨（Mansa Musa，约 1312—1337 年在位）前往麦加朝觐，途经开罗。马克里齐（al-Maqrīzī，1364—1442）也把他到开罗的时间记在 724 年（1324），说他带着大量的礼物和充足的黄金\cite{umari1927}。他的随从里有“一万四千名年轻的女奴”\cite{umari1927}。

十二年后，在马穆鲁克王朝任职的乌马里（Ibn Faḍl Allāh al-ʿUmarī，1301—1349）记下了开罗金价的变化。曼萨·穆萨来的“那一年，一米斯卡勒金的价格没有低于二十五迪拉姆”，而现在“一米斯卡勒不超过二十二迪拉姆，甚至更低。十二年来一直如此”。米斯卡勒（mithqal）是阿拉伯的重量单位，一米斯卡勒金指一米斯卡勒重的金子；迪拉姆是银币。乌马里说，金价跌下来，是因为马里人带到埃及并在那里花掉的黄金数量巨大\cite{umari1927}。以白银计，开罗的金价至少下跌了 12\%，而且十二年没有回升。据乌马里的消息来源说，曼萨·穆萨离开自己的国家时，带了一百驮黄金\cite{umari1927}。

这位国王的黄金来自西非，可在 14 世纪的西非，流通的不只是黄金。盐、铜条、布匹和来自印度洋的贝壳，都在充当货币。两百年后，西班牙人在大西洋的另一边看到了另一番景象。16 世纪的安第斯山上，印加帝国征税、养兵、备荒，几乎不用货币。货币一定要由国家发行吗？一个庞大的帝国能不能完全不用货币？答案要从撒哈拉商路上的黄金、盐与铜说起。

\section{当时的世界}\label{sec:africa-world}

曼萨·穆萨的国家里，钱不止一种。14 世纪的西非人用什么交易、纳税和储藏财富？这些东西又是谁生产的？到了 16 世纪初的安第斯山与中美洲，答案又不一样。

\subsubsection{黄金、盐与铜}
西非的黄金经撒哈拉运往北非与埃及。乌马里关于马里的资料，很多是听伊本·阿米尔·哈吉布（Ibn Amir Hajib）讲的。这位官员在开罗接待过曼萨·穆萨，所以乌马里的记载是转述，不是目击。按这些转述，马里国王不向治下的异教徒征收人头税，而是用他们在矿里开采黄金\cite{umari1927}。曼萨·穆萨还告诉伊本·阿米尔·哈吉布，他的国家有一座铜矿，铜的收入是他最大的税源。铜被运到异教徒黑人的地方出卖，“一百米斯卡勒的这种铜，换六十六又三分之二米斯卡勒的金”\cite{umari1927}。按重量，一份黄金只换一份半的铜。这个比价高得惊人，说明在远离铜矿的地方，铜几乎和黄金一样值钱。

二十多年后，摩洛哥人伊本·白图泰（Ibn Baṭṭūṭa，1304—约 1369）亲自走了这条路。1352 年初，他从西吉尔马萨（Sijilmasa）出发，走了二十五天，到了出产岩盐的塔加扎（Taghaza）\cite{battuta1858}。那里开采盐的是奴隶，一头骆驼通常只能驮两块盐板。他记下了盐价：“一驮这种矿物，在伊瓦拉坦（Iwalatan）卖八到十个米斯卡勒的金，在马里城值二十到三十，有时甚至四十。”\cite{battuta1858}法译本把这里的单位写作“米斯卡勒，或金第纳尔，或杜卡特”。金第纳尔与杜卡特都是金币的名字；在本章的史料里，它们和米斯卡勒金一样，都指一定重量的金。一驮是两块盐板，所以一块盐板在伊瓦拉坦值 4—5 个米斯卡勒，到了马里城值 10—15 个。他还说，那里的黑人用盐作货币\cite{battuta1858}。

回程经过塔凯达（Takedda），他看到了当地的铜矿。铜由男女奴隶冶炼，铸成粗细两种铜条，粗的四百根值一个金杜卡特，细的六百或七百根也值一个。人们用细的买肉和柴，“用粗的买男女奴隶、小米、黄油和小麦”\cite{battuta1858}。粗细两种铜条就是大小两种面额。1353 年他离开塔凯达北返时，同行的商队里有大约六百名女奴\cite{battuta1858}。黄金、盐和铜的开采与运输，处处离不开奴隶；奴隶本身也是这条商路上的一种主要商品。

\subsubsection{布与贝壳}
在马里以东的加奈姆（Kanem），乌马里说，人们以布为货币，每一块布长十肘（肘是长度单位）。当地也用贝壳、玻璃珠、铜块和银币，“但这一切都按布的行市计”\cite{umari1927}。布在这里是\textbf{记账单位}（unit of account），别的东西都用布来报价，即使交易时用的是别的东西。伊本·白图泰从塔凯达到图瓦特（Tuwat）的路上，食物也是用布买的\cite{battuta1858}。

贝壳来自更远的地方。伊本·白图泰在 1344 年离开马尔代夫之前，在那里住过一段时间。他说，当地按数目计贝壳，一万二千枚叫一个 cotta，十万枚叫一个 bostoû。四个 bostoû 换一个金第纳尔，有时便宜到十二个 bostoû 才换一个\cite{battuta1858}。马尔代夫人把贝壳卖给孟加拉人换米，因为贝壳在孟加拉也是通用的货币。也门人买了贝壳，“放在船里当压舱物，代替沙子”\cite{battuta1858}。

几年以后，他在马里城与加奥（Gao）看到同样的贝壳出卖，一千一百五十枚换一个金第纳尔\cite{battuta1858}。在加奥，“居民的买卖用小贝壳，而不用货币”\cite{battuta1858}。乌马里也说，在马里，以贝壳为主要输入品的商人从这项贸易中获利甚丰\cite{umari1927}。

\subsubsection{安第斯山与中美洲}
大西洋对岸的情形不同。西班牙人在 16 世纪征服秘鲁时，印加帝国已经统治着安第斯山的大片地区。加西拉索·德·拉·维加（Garcilaso de la Vega，1539—1616）生于库斯科，父亲是西班牙征服者，母亲是印加王族，晚年在西班牙写成《印卡王室述评》。他说，在他的时代，即他离开秘鲁之前，“以及此后二十年，我的国家都没有铸币。西班牙人买卖时，用马克或盎司称金银的重量”\cite{garcilaso1869}。这里的马克与盎司都是称金银用的重量单位，不是钱币。

印第安人并不缺金银。耶稣会士阿科斯塔（José de Acosta，1540—1600）说，他们的神庙、宫殿与坟墓里有大量金银器物，但“他们不用金银交易或买东西，而是以物易物”\cite{acosta1880}。在中美洲，人们用可可豆交易。阿科斯塔在《印度群岛的自然与道德史》里说，墨西哥各省的人用可可豆代替货币，这种小果实可以买到想要的东西。在秘鲁，人们为同样的目的使用古柯，那是印第安人非常珍视的一种叶子\cite{acosta1880}。

黄金与铜按重量交易，盐与布兼作货币与商品，贝壳从印度洋运来专门充当货币，印加帝国则基本不用货币。同样穿越撒哈拉，为什么黄金能一路运到开罗，盐在两座城之间就差出那么多价钱？

\section{黄金：运费、价差与国王的借款}\label{sec:africa-gold}

曼萨·穆萨只是路过开罗，开罗的金价却十二年没有回升。他的黄金为什么能压低一座大城的金价？穿越撒哈拉的价差由什么决定？他回程时还向商人借了钱。借钱给一位远方的国王，商人要收多高的利息？

\subsubsection{开罗的金价}
乌马里的两个数字，一个是曼萨·穆萨到来那一年金价的下限，一个是十二年后金价的上限，所以 12\% 是跌幅的下限。一米斯卡勒金从至少 25 迪拉姆跌到最多 22 迪拉姆，$1-22/25=12\%$。

马克里齐也记下了金价的下跌。曼萨·穆萨一行在开罗大量购买女奴、歌女与布匹，法译本说，“第纳尔跌到了六迪拉姆”\cite{umari1927}。法译者 Gaudefroy-Demombynes 在注释里认为，第纳尔跌到六迪拉姆不可信，也与乌马里书中的记载相反\cite{umari1927}。问题出在翻译上。这句阿拉伯原文的意思是第纳尔的价格跌了六个迪拉姆，而不是跌到六个迪拉姆。Van Steenbergen 的校勘本就译作“金第纳尔的价格降了六个迪拉姆”\cite{vs2016}。按原文，金价从至少 25 迪拉姆跌了 6 迪拉姆，跌幅至多 24\%（$6/25$），与乌马里的数字（跌幅至少 12\%）并不矛盾。乌马里给出了前后两个价格，本章用他的数字。

一次性带来的黄金，为什么能让金价十二年不回升？这要看开罗与其他地方之间运金的成本。\chref{ch:silver}讲过金银之间的套利，商人在金价低的地方买进黄金，运到金价高的地方卖出。按那里的式~\eqref{eq:silver-arbitrage}，两地的金银比价之差不会超过运费、风险与利息之和。如果开罗的金价长期比别处低 12\% 以上，而把黄金运出埃及的成本低于这个数，商人就会把金子运走，开罗的金价也就回升了。所以，要么黄金运出埃及的成本足够高，要么别处的金价也下跌了，乌马里的记载不能区分这两种情形。曼萨·穆萨的黄金之所以有这么大的影响，还因为它是一次性涌入的大量存货，不是持续的流量。

\subsubsection{运费决定价差}
同一个道理可以推广到任何商品。设一单位商品在 $A$ 地的价格为 $P_A$，在 $B$ 地为 $P_B$。把它从 $A$ 运到 $B$ 的全部成本，包括运费、途中的损耗、利息与风险，合计为 $P_A$ 的比例 $c$。商人在 $A$ 买进、运到 $B$ 卖出，只有当毛收益 $P_B/P_A-1$ 高于成本 $c$ 时才有利可图。运来的货多了，$B$ 地的价格下跌，直到差价不超过成本：
\begin{keyeq}
\frac{P_B}{P_A}-1\;\le\;c .\label{eq:africa-carry}
\end{keyeq}
式~\eqref{eq:silver-arbitrage} 是它在金银之间的特例。反过来，如果一条商路上一直有人在做这门生意，两地的价差就大致等于运输的成本。

伊本·白图泰的盐价可以说明运输的成本有多高。盐从塔加扎运到伊瓦拉坦，再运到马里城，在马里城卖到伊瓦拉坦的 2—3.75 倍（图~\ref{fig:carry_cost}(a)）。两地的价差在 100\% 到 275\% 之间（按“有时四十”算，是 400\%）。按式~\eqref{eq:africa-carry}，这段路的成本至少与价差一样高，最少也相当于盐在伊瓦拉坦的价格。这些成本包括骆驼与赶驼的人、途中的损耗、商人的利息与风险，以及沿途的税。对盐这种又重又便宜的商品，价格的大部分是运费。黄金恰恰相反。一驮黄金的价值远远高于驮运它的成本，所以黄金可以穿越撒哈拉、地中海，一直运到开罗与欧洲。

\fig{carry_cost}{运费限制了两地的价差。(a) 看两根柱子的高低：1352 年一驮盐在马里城卖 20—30 个米斯卡勒的金，有时 40（虚线），是伊瓦拉坦 8—10 个的 2—3.75 倍。这段路的成本至少与盐在伊瓦拉坦的价格相当。(b) 看橙色与褐色柱子的长短：1894 年贝币（当钱用的贝壳）在扎里亚、卡诺比在洛科贾值钱 33\%、50\%，而雇脚夫把一担 95 磅的贝壳运 4 天，工钱已是这担贝壳价值的 48\%，从卡菲一带运到 180 英里外的扎里亚要 96\%。运费大于价差，没有人会把贝壳运过去套利。脚夫的工钱按洛科贾的价格（一英国先令 2,000 枚）折算。数据：伊本·白图泰游记的法译本（1858）；Robinson（1896）。}{fig:carry_cost}

\subsubsection{借给国王的钱}
曼萨·穆萨在开罗花光了钱。乌马里转述伊本·阿米尔·哈吉布的话说，国王回到开罗时已经没有钱了，只好以自己的担保向商人借钱。商人们赚了很多，“以至于三百第纳尔，他们得到七百第纳尔的收益”，后来国王把欠款足额寄还给了他们\cite{umari1927}。马克里齐也说，他不得不向商人大笔借款。回程途中天寒，人与骆驼大量死亡，“与他一同回去的只剩大约三分之一”\cite{umari1927}。

伊本·赫勒敦（Ibn Khaldūn，1332—1406）记下的结局却不同。按他的记载，曼萨·穆萨向开罗的大商人借钱，其中库维克家族（Banū 'l-Kawik）借给他五万第纳尔，他把苏丹赐给他暂住的宫殿卖给了他们。债主派代理人随他去马里收账，代理人死在那里。又派一个人带着儿子去，那人又死在那里。儿子再去，“曼萨·穆萨还没有还钱就死了，他们一文也没有收回”\cite{umari1927}。两份记载说的可能是不同的债主，也可能是同一次借款的两种说法，史料不足以区分。

商人为什么要这么高的利息？借给一位远方国王的钱，债主无法强制收回。国王回到撒哈拉以南，开罗的法庭就管不到他了。设债主是风险中性的，只看期望的收益，国王还款的概率为 $p$。一笔借款如果按时偿还，债主连本带利得到 $1+i$；如果不还，什么也得不到。债主把钱借给国王，而不是投在别的生意上，至少要得到别处的收益 $1+r_f$：
\begin{keyeq}
p\,(1+i)=1+r_f\qquad\Longrightarrow\qquad p=\frac{1+r_f}{1+i} .\label{eq:africa-risk}
\end{keyeq}
利息越高，说明债主估计的还款概率越低；高出 $r_f$ 的部分是对违约风险的补偿。

\begin{example}{开罗的金价与国王的借款}\label{ex:africa-gold}
（1）乌马里说金价从至少 25 迪拉姆跌到最多 22 迪拉姆，跌幅至少是多少？（2）“三百第纳尔得到七百第纳尔的收益”有两种读法：连本带利一共还 700，或者利润 700、连本带利还 1,000。两种读法下，这笔借款的利率 $i$ 各是多少？（3）设债主的钱投在别处，在同样的期限里可以赚 $r_f=15\%$（教学假设；\chref{ch:islamic}表~\ref{tab:islamic-rates} 中意大利商业贷款的年息为 10\%—18\%）。按式~\eqref{eq:africa-risk}，两种读法下债主估计的还款概率各是多少？如果 $r_f=0$ 呢？
\solution （1）$1-22/25=12\%$；这是下限，因为两个数字一个是下限、一个是上限。（2）还 700：$i=700/300-1\approx133.3\%$；还 1,000：$i=1\,000/300-1\approx233.3\%$。（3）$p=1.15/2.333\approx0.49$，$p=1.15/3.333\approx0.345$；$r_f=0$ 时，$p=1/2.333\approx0.43$，$p=1/3.333=0.30$。借款的期限史料没有说。他是朝觐归来、回到开罗时借的钱，钱要等他回到马里以后再送回开罗，一来一回要走好几个月，期限可能接近一年，也可能更长。期限短于一年，同期的 $r_f$ 低于 15\%，这里的 $p$ 偏高；长于一年，$r_f$ 高于 15\%，$p$ 偏低。债主若厌恶风险，要求的补偿更多，同样的利率对应的 $p$ 会更高。这些偏差改变不了数量级：无论哪种读法，商人都把这位带着一百驮黄金来朝觐的国王看成一个很可能赖账的借款人。事后看，按乌马里的消息来源，借款还了；按伊本·赫勒敦，库维克家族的钱没有收回。
\end{example}

借给远方君主的钱要靠什么收回，是主权债务的老问题。\chref{ch:silver}的腓力二世一再停付，仍然借得到钱，因为他的热那亚银行家可以集体停贷，而他离不开下一笔借款。曼萨·穆萨的债主没有这样的手段，也不指望他再来开罗借钱。

\section{贝币：从马尔代夫到西非}\label{sec:africa-cowrie}

同一枚贝壳，在马里比在马尔代夫值钱几百倍。西非人把这种来自印度洋的小海贝当钱用，这就是贝币。它怎样成了西非的货币？它的价值由什么决定，为什么在 19 世纪一路贬值？统治者又能对它做什么？

\subsubsection{贝币的单位}
贝币是成千上万地使用的。西非海岸的通行办法是把贝壳穿成串。按达尔泽尔（Archibald Dalzel，1740—1818）《达荷美史》（1793）的注释，“散的四十枚贝壳为一串；穿起来的为 39 枚，第 40 枚抵作穿孔与穿线的工钱”\cite{dalzel1793}。

串之上还有更大的单位。达荷美人的大单位叫“头”（cabess，葡萄牙文 cabeça）。该书编者在注释（署名 J. F.）里列了一张表：40 枚为一串，200 枚为一个 galhina，1,000 枚为一个 ackey，4,000 枚为一头，16,000 枚为一“贸易盎司”（ounce trade）\cite{dalzel1793}。

贸易盎司是一个记账单位，和加西拉索说的称金银的盎司不是一回事。J. F. 说，这里的盎司指的是价值。盎司有两种，金盎司值 4 英镑，贸易盎司只值它的一半。英镑是英国的货币，一英镑合 20 先令、240 便士。所以 16,000 枚值 40 先令，一头 4,000 枚值 10 先令。书中另一处注释也说，一头贝币是十先令。将近一百年后，埃利斯（A. B. Ellis，1852—1894）记下的一头是五十串，即 2,000 枚\cite{ellis1890}。各地各时的“头”大小不一，下面用到时都写明枚数。

贝币很重。J. F. 说，一贸易盎司的贝壳重约四十二磅（常衡）。磅（pound）是英国的重量单位，与货币英镑同名，一磅（常衡）为 453.59 克，一长吨为 2,240 磅。1772 年，英国奴隶贩子诺里斯（Robert Norris，卒于 1791）在达荷美国王的宫廷里看到，国王赏给乐师“八十盎司贝币，由他的三百二十个妻子送来，每人用铜盘端着一头，即四千枚”\cite{dalzel1793}。$320\times4\,000=1\,280\,000$ 枚，正是 80 个贸易盎司，值 160 英镑。按 J. F. 的重量，共重 3,360 磅，即 1.5 长吨，与原注说的“约一吨半”相符，每人端着 10.5 磅。

\subsubsection{一万里的价差}
伊本·白图泰在马尔代夫与马里看到的价格，相差得惊人。在马尔代夫，一个金第纳尔买 400,000 枚，便宜时买 1,200,000 枚；在马里，只买 1,150 枚。一枚贝壳在马里的价值，是在马尔代夫的 348—1,043 倍（图~\ref{fig:cowrie_value}(a)）。

\begin{example}{贝壳从马尔代夫到马里}\label{ex:africa-cowrie}
（1）按伊本·白图泰的数字，一个金第纳尔在马尔代夫买进的贝壳运到马里，能卖得几个第纳尔？（2）按式~\eqref{eq:africa-carry}，这条路的成本至少要多高，马尔代夫与马里之间的价差才能维持？（3）为什么贝壳的价差可以这么大，而黄金的价差不能？
\solution （1）$400\,000/1\,150\approx348$ 个第纳尔；便宜时 $1\,200\,000/1\,150\approx1\,043$ 个。（2）价差 $P_B/P_A-1$ 至少为 347，即 34,700\%；只要商路上一直有人运贝壳，全部成本就大致是这么多。（3）贝壳在产地几乎一文不值，运输、转手与风险的成本加起来，比它在产地的价格高几百倍。海上的一段路几乎不花钱，因为贝壳可以当压舱物；花钱的是之后的转运，尤其是陆路上的骆驼商队。黄金正好相反，一驮的价值远高于运费，$c$ 很小，价差也就很小。贝壳的价差反映的不是马里人“不懂行情”，而是把它运过来有多难。
\end{example}

到了大西洋贸易的时代，贝币多由欧洲商人经海路运到西非海岸，用来购买非洲的出口品，其中包括奴隶。Gardner 说，到 19 世纪初，“仅英国人每年就运往非洲大约 100 吨贝壳”\cite{gardner2015}。

\subsubsection{运不动的钱}
贝币便于小额交易，到了大额交易就成了负担。英国传教士罗宾逊（Charles Henry Robinson，1861—1925）在 1894—1895 年穿越豪萨诸邦，留下了一份精确的账\cite{robinson1896}。在洛科（Loko），他雇一个脚夫，付“一万枚贝币，按英国的钱是五先令，把一担九十五磅的货物运四天”。在卡菲（Kaffi），他卖布凑出二十二万枚贝币，“约合英国的钱六镑，重约一千磅”\cite{robinson1896}。一磅约 220 枚，一枚约 2.06 克。按洛科贾（Lokoja）的价格，即两千枚一先令，一担 95 磅的贝壳只有 20,900 枚，值 10.45 先令。

\begin{example}{把贝币运到扎里亚}\label{ex:africa-carry}
罗宾逊说，两千枚贝币“在洛科贾值一先令，在扎里亚值一先令四便士，在卡诺值一先令六便士”；离开卡菲以后，他付给每个脚夫“两万枚贝币，约十先令”，把一担货运到 180 英里外的扎里亚\cite{robinson1896}。（1）雇脚夫把一担 95 磅的贝壳运 4 天，工钱占这担贝壳价值的多少？运到扎里亚呢？（2）贝币在扎里亚、卡诺比在洛科贾值钱多少？（3）按式~\eqref{eq:africa-carry}，商人会不会把贝壳从洛科贾运到扎里亚套利？
\solution （1）4 天：$10\,000/20\,900\approx47.8\%$，平均每天约 12\%；到扎里亚：$20\,000/20\,900\approx95.7\%$。（2）扎里亚：$16/12-1\approx33.3\%$；卡诺：$18/12-1=50\%$。（3）不会：从卡菲一带运到扎里亚的工钱已经几乎吃掉整担贝壳，从洛科贾出发路还要更远，而价差只有三分之一。贝币在各个市场之间的价差可以长期存在，因为运它的成本比价差还大（图~\ref{fig:carry_cost}(b)）。
\end{example}

运不动的钱带来几种后果。第一，大额的东西很难卖掉。罗宾逊的马病倒了，却卖不出去，因为“它的价值换成贝币，要多雇十五个脚夫来背，而我们得把他们背的钱全付给他们，还不止”\cite{robinson1896}。按每担 95 磅算，十五担是 313,500 枚，按洛科贾的价格约值 7.8 镑。他只好把马丢下。

第二，凑钱要花时间。在卡菲，他们“足足耽搁了十天，因为很难凑够发给脚夫当伙食钱的贝币”。那里没有商店，也没有批发商，只好雇人把布拿到市场上零卖\cite{robinson1896}。这 22 万枚贝币重约一千磅，要 10.5 个脚夫来背。在扎里亚，他们要凑将近七十五万枚贝币付给脚夫，随身带的易货商品“除非亏得离谱，不可能很快卖掉”\cite{robinson1896}。

第三，便于携带的钱有溢价。玛丽亚·特蕾莎银元是一种银币。罗宾逊说，它在维也纳铸造，上面的年份是 1780 年，专门为出口到中非而造，“含银约值两先令，在苏丹的购买力从三先令到四先令六便士不等”\cite{robinson1896}。这里的“苏丹”指撒哈拉以南的草原地带，不是今天的苏丹国。银元的价值比它的含银高出 50\%—125\%，高出的部分叫溢价。在卡菲，一枚银元换到约三先令六便士的贝币，高出 75\%。一枚银元可以放进口袋，同样价值的贝币要几千枚。这份溢价就是\textbf{便携性}的价格。但银元的市场很浅。他们在扎里亚按每枚五千贝币卖了大约一百四十枚玛丽亚·特蕾莎银元，此后几天，银元就贬值了约 20\%，即从五千跌到四千\cite{robinson1896}。140 枚银元换来 700,000 枚贝币，正好补上他们要凑的那笔钱的大头。

\fig{cowrie_value}{贝币的价值随距离与时间变化。(a) 看柱子的长短（对数横轴）：1344 年在马尔代夫，一个金第纳尔买 40 万到 120 万枚贝壳，1352 年在马里与加奥只买 1,150 枚。1851 年在卡诺，一米斯卡勒金买 4,000 枚，廷巴克图的名义价格为 3,500—4,500 枚。(b) 看圆点的走向（对数纵轴）：一先令换的贝壳，1720 年为 200 枚，1793 年按 J. F. 的记账单位为 400 枚，1863 年维达约 1,111 枚、拉各斯约 1,511 枚，1894 年洛科贾 2,000 枚。1720—1894 年平均每年多换 1.33\%。各点的地点不同，1863 年按每“头”2,000 枚折算。数据：伊本·白图泰游记的法译本（1858）；Barth（1857）；Gardner（2015）；Dalzel（1793）；Burton（1864）；Ellis（1890）；Robinson（1896）。}{fig:cowrie_value}

第四，最大的交易干脆不用贝币。罗宾逊在扎里亚的市场上一次看到多达三百个待卖的奴隶，一个奴隶卖“十万到三十万枚贝币，即三到九英镑”\cite{robinson1896}。按一磅 220 枚算，一个奴隶的价钱是 455—1,364 磅贝壳，要 4.8—14.4 个脚夫来背。于是人本身成了大额支付的手段。罗宾逊说，奴隶有两种用途，“一是作为货币，用在国家贡赋这样的大额支付上，二是作为脚夫”。豪萨各邦向索科托（Sokoto）缴纳的贡赋，至少四分之三是用奴隶缴纳的。据说阿达马瓦（Adamawa）的国王每年送一万名奴隶，罗宾逊后来听到的另一个说法是两千名。卡诺送得最少，据说约一百名，贡赋的大部分用布缴纳\cite{robinson1896}。贝币运不动，大额的支付就落到了被奴役的人身上。

\begin{intuition}[运不动的钱留在原地]
一种货币能在多大范围内通用，取决于把它运来运去要花多少钱。黄金每单位重量的价值很高，可以在整个旧大陆之间流动，各地的金价差不多；贝壳每单位重量的价值很低，各个市场的价格可以差出三成、五成，到了印度洋与西非之间，差出几百倍。所以贝币天生是“本地的”货币。在一个市场之内，它是方便的零钱；在市场之间，它的价值由运费隔开。大额的支付与远距离的价值储藏，要用别的东西：黄金、银元、布，乃至奴隶。
\end{intuition}

\subsubsection{统治者与贝币}
贝币不是国家发行的，但统治者并不只是旁观。1851 年，德国旅行家巴特（Heinrich Barth，1821—1865）在卡诺看到，贝币是市场的日常货币，2,500 枚等于一个西班牙或奥地利银元。以金计，一米斯卡勒通常维持在四千枚贝壳，巴特说这是一个普遍的标准\cite{barth1857}。按这两个数字，一米斯卡勒金值 1.6 个银元。他也注意到贝币的不便。在内陆，贝壳不像海岸附近那样穿成串，“必须一枚一枚地数”。城里的官员习惯把它们装进每袋两万枚的草袋，私人收下时仍要重新数。他的一位商人朋友在五六个人的帮助下数完五十万枚，即 25 袋，巴特称之为“真正英雄般的工作”\cite{barth1857}。

往东的博尔努（Bornu）情形不同。巴特说，那里原来以棉布条为货币，近来才开始被贝币取代，而贝币是被引进来的，“与其说是出于居民的自然需要，不如说是出于统治者的投机”\cite{barth1857}。博尔努的记账单位叫罗特尔（rotl），原来是一磅铜。铜磅早已不用，名字却留了下来，32 枚贝壳算一罗特尔。

银元与罗特尔的比价波动极大。巴特说，这至少部分来自当权者的投机，主要是他的朋友哈吉·贝希尔（Hajj Beshir）的投机。他承认无法为朋友辩护，因为这种投机“大大损害了公众”。每当贝希尔收集了大批贝币、想让它流通，“银元会突然跌到四十五或五十罗特尔，而在别的时候，它能换到一百罗特尔，即三千二百枚贝壳，比在卡诺多七百枚”\cite{barth1857}。一银元的价格在 1,440 枚到 3,200 枚之间变动，高低相差 2.22 倍。卡诺则不同。在巴特看来，那里只有一种标准的钱币，只要不突然把太多的银元投入流通，它的价值总是不变\cite{barth1857}。

几种货币并存而没有固定的比价，交易就会退回到一连串的交换。在博尔努的都城库卡瓦（Kukawa），卖粮的小农不收贝币，也很少收银元。想买粮的人“必须先用银元换贝币，或者说买贝币”，再用贝币买一件衬衫，经过一番讨价还价，才能买到粮食\cite{barth1857}。

巴特在脚注里还说，贝币大量流入博尔努，抽空了豪萨地区的贝币，影响了 1849 年以后约鲁巴与海岸对贝币的需求\cite{barth1857}。他在卡诺也注意到，1848 年以来，海岸上对这种贝壳的需求大大增加了。

统治者对贝币动手最彻底的例子，是达荷美的国王。

\begin{case}{维达 · 1774—1789 · 阿达洪祖二世的垄断与短串}
达荷美（Dahomey）位于今天的贝宁，经由沿海的维达（Whydah，今乌伊达）港与欧洲商人做奴隶生意。按达尔泽尔的记载，国王博萨·阿哈迪（Bossa Ahadee，即 Tegbesu，1740—1774 年在位）约七十岁时，于 1774 年 5 月 17 日去世，继位的是阿达洪祖二世（Adahoonzou II）\cite{dalzel1793}，即 Kpengla（1774—1789 年在位）。达尔泽尔曾任英国在维达的商馆总督。据书中说，阿达洪祖二世的事迹采自当时的英国维达总督阿布森（Lionel Abson）的报告。

阿达洪祖二世打仗不顺，为了弥补战争的失利，他决心试一试贸易垄断的效果。他下令，任何商人在任何市场上，“买一个男奴不得付出三十二头以上的贝币，买一个女奴不得超过二十六头”。他派人监视商人，凡违令的，“没收他们的奴隶，连他们本人一起，如果能卖掉的话；如果犯令的商人卖不掉，就砍头”\cite{dalzel1793}。

按每头 10 先令，男奴的限价是 16 英镑，女奴是 13 英镑。而 1786 年前后，维达的港口里停着六艘奴隶船，货源奇缺，一个壮年奴隶的价格将近三十英镑\cite{dalzel1793}。限价只有市价的 53.3\%（男）与 43.3\%（女）。两个数字的年份不完全相同，只能看出大致的差距。

国王还要分一杯羹。每个商人“有两个奴隶的，要把其中一个、而且是最好的一个交给他；有三个的，两个归国王”，国王按自己定的价格付给穿成串的贝币\cite{dalzel1793}。穿串的是国王宫中的妇女，她们“特意让每串少三到六枚”\cite{dalzel1793}。

一串本应穿 39 枚，少 3—6 枚，就是少付 7.7\%—15.4\%。按\chref{ch:medieval}的式~\eqref{eq:medieval-seigniorage}，这等于把每串的“含量”从 39 枚降到 36 或 33 枚，$MP/ME$ 只有 92.3\% 或 84.6\%。短串是一种铸币税，付款的一方暗中减少了每个单位所含的东西。一个男奴按限价付 16 英镑，用短串付，实际只付 13.5—14.8 英镑，是 1786 年市价的 45\%—49\%。

国王的说法是，这些措施是为了“让财富留在本国”\cite{dalzel1793}。结果是商路改道，一边流向维米（Wemey），去波多诺伏（Porto Novo）的市场，另一边流向佩希（Peshie），去波波（Popo）与基塔（Quitta）的市场。达尔泽尔写道，“许多人破产，一些拥有六七十个奴隶的财产的大商人沦为乞丐”。国王又加征关税，一时让收入显得很兴旺，而维达好几年完全没有贸易\cite{dalzel1793}。1789 年 4 月 17 日，阿达洪祖二世死于天花\cite{dalzel1793}。

这段记载要放在它的来源里读。达尔泽尔与阿布森都是奴隶贸易的当事人，书的编者在序言里为奴隶贸易辩护，说“我们确信，这对不幸的勇士是一种仁慈”\cite{dalzel1793}。达尔泽尔称阿达洪祖二世为“人类的祸害”\cite{dalzel1793}，但他记下的这些罪状，损害的首先是奴隶贩子的生意。而这门生意交易的“商品”，是人。
\end{case}

\subsubsection{19 世纪的贝币通胀}
贝币的长期价值可以从几个相隔很远的数字看出来（图~\ref{fig:cowrie_value}(b)）。Gardner 说，1720 年，一个 100 串（即 4,000 枚）的“大头”值大约 1 英镑\cite{gardner2015}。这相当于一先令合 200 枚。J. F. 在 1793 年把一头 4,000 枚记作 10 先令，一先令合 400 枚。1863 年，伯顿（Richard F. Burton，1821—1890）出使达荷美国王格莱莱（Gelele，1858—1889 年在位）。他说，“一银元（四先令六便士）现在在维达与阿波美买 2½ 头，在拉各斯与阿贝奥库塔买 3 头又 20 串”\cite{burton1864}。按每头 2,000 枚，一先令在维达合约 1,111 枚，在拉各斯合约 1,511 枚。伯顿自己算的一头值一先令九便士半，与 $54/2.5=21.6$ 便士只差一点。1894 年，罗宾逊在洛科贾按两千枚一先令计算。

从 1720 年到 1894 年，一先令换的贝壳增加到 10 倍。这 174 年里，平均每年多换 1.33\%。从 1793 年到 1863 年，它增加到 2.78 倍，平均每年 1.47\%。英镑在这一时期大体钉住黄金，一英镑合多少黄金大致不变，所以这也是贝币对黄金的贬值。年均一个多百分点，以今天的标准是温和的通胀，但它持续了将近两个世纪。更长的时段里，贬值慢得多。1352 年在马里，一个金第纳尔买 1,150 枚；1851 年在卡诺，一米斯卡勒金买 4,000 枚。499 年里增加到 3.48 倍，平均每年只有 0.25\%。

贬值的原因是供给。伯顿说，“必须记住，贝币是商品，价格随行就市”；他提到桑给巴尔的“蓝贝”，即大而便宜的贝壳\cite{burton1864}。Gardner 说，19 世纪欧洲对非洲产品的需求增长，商人运来的贝壳越来越多。在马尔代夫的贝壳之外，又加上了桑给巴尔出产的一种更大、更多的贝壳（Cypraea annulus）\cite{gardner2015}。

在 Gardner 看来，19 世纪的“贝币通胀”增加了使用贝币作交易媒介的成本，包括搬运、清点与储存的成本。两份重量可以看出这种变化。J. F. 的一贸易盎司是 16,000 枚，重 42 磅，算下来一磅约 381 枚，一枚约 1.19 克。罗宾逊的贝壳一磅只有约 220 枚，一枚约 2.06 克。按 J. F. 的记账，1793 年一磅贝壳约值 0.95 先令。按罗宾逊的数字，1894 年只值 0.12 先令，相差 7.9 倍。贝壳变大了，每磅的价值却低了，运它的成本相对于它的价值就高了。

\begin{compare}{商品货币与国家货币}{贝币、盐、布、可可豆}{铸币与纸币}
\cmprow{价值从哪里来}{运来的成本与用途：压舱物、装饰、可以吃喝}{国家的接受与法令：纳税、法偿}
\cmprow{数量由谁决定}{远方的供给者与运费：马尔代夫、桑给巴尔、欧洲商人}{发行者：铸币厂与发行机构}
\cmprow{贬值从哪里来}{供给增加，运费下降：19 世纪的“贝币通胀”}{发行过多、减少含量（第~\ref{ch:medieval}、\ref{ch:songming}~章）}
\cmprow{统治者能做什么}{限价、短串、垄断（达荷美）；囤积投机（博尔努）；规定用什么纳税（索科托收奴隶）}{收取铸币税；规定法偿；禁止外币}
\cmprow{主要的缺点}{笨重、难数，各个市场的价差大}{取决于发行者能不能自我约束}
\cmpsum{贝币不是国家发行的，但统治者可以通过限价、短串、投机与规定用什么纳税来左右它；殖民政府用禁令与新货币取代它，也花了几十年。}
\end{compare}

\section{没有货币的帝国：印加}\label{sec:africa-inca}

印加帝国的百姓纳税，交的不是钱，也不是自家的收成，而是劳动。一个横跨安第斯山的帝国，怎样在不用货币的情况下征税、养兵、备荒？这种安排靠什么维持，又有什么限制？

\subsubsection{只收劳役的税}
按加西拉索的说法，太阳（即神庙）与印加国王各有自己的土地，收成分开存放。百姓耕种这些土地，种子由土地的主人提供，干活的人也由它供养，“所以印第安人只需献出个人的劳动作为贡赋。臣民不必从自己的收成里向印加缴纳任何东西”\cite{garcilaso1869}。

加西拉索还引用了耶稣会士布拉斯·瓦莱拉（Blas Valera，1545—1597）的记载。按瓦莱拉的说法，印加的贡赋如此适度，古代的君主里，连罗马的奥古斯都与庇护在内，没有一位能在这一点上与印加诸王相比。许多印第安人缴的，几乎不值四个里亚尔（西班牙的一种钱币）。所有男孩直到二十五岁免役，五十岁以上的男子也免除贡赋\cite{garcilaso1869}。瓦莱拉列出的第三条法律是，任何人都不必从自己的财产里拿出东西来纳贡，只需缴纳劳动、时间或手艺，用来为国王或国家服务\cite{garcilaso1869}。按这些记载，一个男子从 25 岁到 50 岁，有 25 年要为国家出劳役。“他们说一个人富有，是指他有子女”\cite{garcilaso1869}。在一个按人头出劳役的社会里，财富就是人手。

劳役税不需要货币。国家不需要先把百姓的劳动变成钱，再用钱去买粮食、布与道路，而是直接让百姓去种地、织布、修路。

\subsubsection{仓库}
劳役生产出来的东西要储存起来。加西拉索说，有三种仓库存放收成和其他贡赋，其中一种存放的是为百姓储备的粮食，以防荒年。“王家大道上还有别的仓库，每隔三里格一座”\cite{garcilaso1869}。里格（league）是西班牙的长度单位。他还特别说到冻干的马铃薯（chuño）：“这样处理过的马铃薯可以保存很久”\cite{garcilaso1869}。

仓库的账可以写成一个简单的方程。设 $S_t$ 为第 $t$ 年初的库存，每年由劳役交进 $I$，平时的开支（军队、官吏、工程）为 $G$，库存每年损耗比例 $\delta$。下一年的库存为
\begin{keyeq}
S_{t+1}=(1-\delta)\,S_t+I-G_t,\qquad S^{*}=\frac{I-G}{\delta} .\label{eq:africa-store}
\end{keyeq}
在平常年份，库存会趋向一个稳态 $S^{*}$。库存越多，每年损耗掉的越多，直到损耗正好等于每年的结余 $I-G$。$\delta$ 越小，能存下的越多。冻干马铃薯的用处就在这里，它让 $\delta$ 变小，仓库才能存得住一年以上的粮食。

\begin{example}{仓库能存多少粮}\label{ex:africa-store}
以一年的消费为 100。设每年交进的减去平常的开支，结余为 5，即一年消费的 5\%（教学假设，下同）。耐储存的冻干食物每年损耗 $\delta=5\%$，不耐储存的每年损耗 25\%。（1）两种情形下，稳态库存各是多少？够吃几个月？（2）库存与稳态的差距每年缩小 $\delta$，差距减半要几年？（3）设仓库从零开始建，第 30 年遇上饥荒，要额外动用 40。两种情形下，仓库够不够？
\solution （1）由式~\eqref{eq:africa-store}，$S^{*}=5/0.05=100$，够吃 12 个月；$S^{*}=5/0.25=20$，只够 2.4 个月。（2）差距每年乘以 $1-\delta$，减半要 $\ln2/[-\ln(1-\delta)]$ 年：$\delta=5\%$ 时约 13.5 年，$\delta=25\%$ 时约 2.4 年。（3）第 30 年饥荒之前，前者的库存为 $100\times(1-0.95^{30})\approx78.5$，动用 40 之后还剩 38.5；后者早已到了稳态 20，缺 20（图~\ref{fig:inca_store}）。损耗率决定了仓储能提供多少保险。这几个数字只说明这个关系，不是印加仓库的实际规模。
\end{example}

\fig{inca_store}{仓库能存多少，取决于东西坏得多快。(a) 看两条线的高度：每年结余为一年消费的 5\%，不易坏的食物（每年损耗 5\%）的库存趋向 100，第 30 年饥荒多用 40 之后还剩 38.5；易坏的食物（每年损耗 25\%）的库存只有 20，饥荒时缺 20。(b) 看曲线的形状：稳态库存够吃的月数 $12\times5\%/\delta$ 随损耗率迅速下降，5\% 时 12 个月，25\% 时 2.4 个月。教学模型，式~\eqref{eq:africa-store}；参数是本书的假设。}{fig:inca_store}

\subsubsection{结绳：没有货币的账}
没有货币，国家就要自己记住每户出了多少劳役、仓库里有什么、谁领了什么。印加人用结绳记账，这些绳结叫奇普（quipu）。加西拉索说，奇普由称为奇普卡马约（Quipu-camayu）的印第安人掌管。每个村子“有四个记账的人，人口多的地方，从这个数目到二十或三十个”，国王希望每个村子都有几个，互相核对\cite{garcilaso1869}。按瓦莱拉的说法，分摊贡赋时，他们用绳结，也用小石子来计算\cite{garcilaso1869}。至于百姓之间，加西拉索说，他们不懂租赁与买卖，“因为没有货币。他们用一种食物交换另一种食物，仅此而已”\cite{garcilaso1869}。

\begin{intuition}[货币是一种记忆]
为什么一个社会需要货币？Kocherlakota（1998）的回答是，货币是一种记账的工具。他把“记忆”定义为一个人对他过去直接或间接接触过的所有人的完整历史的了解，并证明了“在有货币的环境中可行的任何配置，在同一环境中有记忆时也可行”。在这个意义上，“货币等价于一种原始形式的记忆”\cite{kocherlakota1998}。手里的钱记录的是，我过去为别人提供过东西，所以现在有权从别人那里拿东西。如果有一本记下每个人全部历史的账，就不需要这种凭证。印加的结绳就是这样一本账，只是它只记录每户与国家之间的义务，不记录百姓之间的往来。代价是记忆很贵。每个村子要有 4 到 30 个专门的记账人，还要互相核对。货币把这本账分散地记在每个人的口袋里，谁也不必知道别人的历史。\chref{ch:mesopotamia}的神庙与宫廷也是靠记账而不是货币来组织生产的。
\end{intuition}

\begin{compare}{劳役税与货币税}{劳役税（印加）}{货币税（一条鞭法以后的明朝）}
\cmprow{国家收到什么}{劳动，产品由国家自己保管、运输}{白银，可以在市场上买需要的东西}
\cmprow{需要什么样的账}{每户每人的劳役都要记录，每村 4—30 个记账人}{只记收了多少钱，谁交的钱都一样}
\cmprow{储存的成本}{仓库里的东西每年损耗 $\delta$}{金属几乎不坏，但购买力随银价变动}
\cmprow{收成与价格的风险由谁承担}{按理想化的记载，歉收由国家的仓库救济}{纳税人：税额以银计，歉收也要照交；粮价下跌时实际税负加重（\chref{ch:silver}）}
\cmprow{前提条件}{国家能直接组织与监督劳动}{有市场，也有足够的货币供给}
\cmpsum{劳役税让国家绕开了货币，代价是国家必须自己记账、保管和运输；货币税把这些事交给了市场，也把风险交给了纳税人。}
\end{compare}

\subsubsection{殖民时期的米塔制}
西班牙人保留了劳役和它的名字，却改变了它的用途。1573—1812 年，秘鲁与玻利维亚实行“一种大规模的强迫采矿劳役制度”，即米塔制（mita）\cite{dell2010}，把指定地区的印第安人征派到波托西等地的矿山。阿科斯塔描写了波托西的矿井：“他们在这些矿里劳动，终日处在黑暗之中”\cite{acosta1880}。矿工沿着梯子往上爬，一个人通常背两阿罗巴的重量。阿罗巴（arroba）与英担都是重量单位，阿科斯塔说四阿罗巴就是一英担，所以两阿罗巴是半英担。矿井从底到口，通常有一百五十多个人高\cite{acosta1880}。米塔制前后实行了 239 年。

矿山把货币经济带进了安第斯山。阿科斯塔说，波托西的古柯叶生意“每年超过五十万比索”（英译作 dollars）\cite{acosta1880}。比索是西班牙的货币单位。白银从这里运往塞维利亚与马尼拉，就是\chref{ch:silver}讲的那股流动。

米塔制的影响延续到今天。Dell（2010）用断点回归比较米塔制边界两侧的地区，发现米塔制使家庭消费降低约 25\%，使受影响地区今天儿童发育迟缓的比例提高约 6 个百分点。按 Dell 的分析，这种影响是通过米塔制对土地制度与公共物品提供的影响延续下来的\cite{dell2010}。

\begin{compare}{印加的劳役与殖民时期的米塔}{印加的劳役}{殖民时期的米塔（1573—1812）}
\cmprow{劳动做什么}{耕种太阳与国王的土地，为国王或国家服务}{到波托西等地的银矿采银}
\cmprow{谁要服役}{按瓦莱拉的记载，25 岁到 50 岁的男子}{指定地区的印第安人}
\cmprow{产出归谁}{进入国家的仓库，一部分用来备荒}{白银运往塞维利亚与马尼拉}
\cmprow{史料}{加西拉索与瓦莱拉在征服之后的追述，有美化}{阿科斯塔的描写；Dell 的计量研究}
\cmprow{长期后果}{史料不足以判断}{今天的家庭消费低约 25\%，儿童发育迟缓高约 6 个百分点}
\cmpsum{名字相同，前者是以劳动代替货币的税，后者把这种税改造成了为白银生产服务的强迫劳动。}
\end{compare}

\section{可可豆：会坏的货币}\label{sec:africa-cacao}

西班牙人在中美洲看到，当地人把一种可以喝掉的东西当钱用。可可豆会受潮、会被虫蛀，也会被人喝掉。这样的东西能不能当货币？持有这样的货币意味着什么？

\subsubsection{中美洲的可可豆}
西班牙编年史家奥维多（Gonzalo Fernández de Oviedo，1478—1557）在尼加拉瓜看到，当地人很看重可可豆。他们给它的估价，就像基督徒与其他民族看待黄金与货币一样。他把可可豆叫作杏仁，记下了价格：“一只兔子值十粒这种杏仁，四粒杏仁可以换八个果子”；“一个奴隶值一百粒左右这种杏仁”\cite{oviedo1851}。一个人的价钱，是十只兔子。

可可豆有货币的种种特征。有人作假，奥维多说，就连这些杏仁里也有骗人的花样；也有人用这种货币或杏仁交纳贡赋\cite{oviedo1851}。它还有一个金属货币没有的特点，就是可以喝掉。奥维多说，把可可豆做成饮料喝掉，“无非是故意让自己变穷，把钱吞下去，或者扔到会丢掉的地方”\cite{oviedo1851}。

阿科斯塔说，在墨西哥，五粒可可豆买一样东西，三十粒买另一样，人们还用它施舍乞丐。可可豆是新西班牙最大的买卖之一，“因为它是干果，可以久存而不坏”。前一年，一个英国海盗在瓜图尔科（Guatulco）港烧掉了十万驮以上的可可豆\cite{acosta1880}。

意大利人文学者彼得·马蒂尔（Pietro Martire d'Anghiera，1457—1526）的看法正好相反。他说：“啊，有福的货币！它不仅给人类一种有益而可口的饮料，还使它的持有者不至于陷入地狱般的贪婪，因为它不能堆积，也不能长久囤积。”他也说，有些酋长只用这种果实纳税\cite{martyr1912}。阿科斯塔说它能久存，马蒂尔说它不能长久囤积。两人都没有给出数字，可能都对，只是比较的对象不同。比起鲜果，它能放很久；比起黄金，它终究会坏。

\subsubsection{持有会坏的货币}
设一份可可豆每年损耗比例 $\delta$（受潮、虫蛀、被喝掉），以可可豆计的物价从 $P_t$ 变为 $P_{t+1}$。年初持有一份可可豆，年底剩下 $1-\delta$ 份，每份能买的东西是原来的 $P_t/P_{t+1}$。持有它的实际收益率 $r^{m}$ 满足
\begin{keyeq}
1+r^{m}=(1-\delta)\,\frac{P_t}{P_{t+1}} .\label{eq:africa-hold}
\end{keyeq}
金属货币的 $\delta$ 接近零，持有它的收益只取决于物价；可可豆的 $\delta$ 大于零，即使物价不变，持有它的收益也是负的。马蒂尔说的“不能囤积”，用今天的话说，就是这种货币的\textbf{持有收益}为负。人们不愿意把财富长期放在里面，会尽快把它花掉，或者换成别的东西储藏。

\begin{example}{持有可可豆的收益}\label{ex:africa-cacao}
设可可豆每年损耗 10\%（教学假设）。（1）以可可豆计的物价不变，持有一年的实际收益率是多少？（2）丰收之年可可豆多了，以可可豆计的物价上涨 5\%（教学假设），收益率是多少？（3）1720—1894 年，一先令换的贝壳平均每年多 1.33\%。贝壳几乎不坏（$\delta\approx0$），按式~\eqref{eq:africa-hold}，以黄金计，持有贝币的年收益率约是多少？
\solution （1）$r^{m}=0.9-1=-10\%$。（2）$0.9/1.05-1\approx-14.3\%$：损耗与通胀叠加在一起。（3）以黄金计，贝币的“物价”每年上涨 1.33\%，$1/1.0133-1\approx-1.31\%$。贝币的问题不在于会坏，而在于供给会增加；可可豆两样都有。持有收益为负的货币适合用来交易，不适合用来储藏财富。
\end{example}

可可豆还有另一面，它的供给取决于收成。丰收之年货币多，歉收之年货币少，而且人们还会把它喝掉。这与\chref{ch:silver}的白银相似，都是一种由自然与运输决定供给的货币，只是可可豆的供给变化要快得多。

\section{数据检验与反例}\label{sec:africa-test}

曼萨·穆萨的国家里同时流通着好几种钱，印加帝国却几乎不用钱。从这两头出发，前面几节得出了三个说法。价差受运费约束（\secref{sec:africa-gold}）。贝币不由国家发行，统治者却可以左右它（\secref{sec:africa-cowrie}）。印加帝国以劳役与仓储代替货币（\secref{sec:africa-inca}）。证据在哪些地方支持它们，在哪些地方不足？

\subsubsection{货币一定要由国家发行吗}
证据大多说“不必”。贝币是商人从印度洋运来的，乌马里说贩运贝币的商人获利甚丰。它的价值由运费与供给决定，各个市场的价格可以相差三成、五成，到了印度洋与西非之间相差几百倍。盐、布、铜条与可可豆也都不是国家发行的。

但“不由国家发行”不等于“与国家无关”。加奈姆以布为记账单位。索科托的贡赋至少四分之三收奴隶，卡诺的贡赋大部分收布，尼加拉瓜也有人用可可豆交纳贡赋。国家规定用什么纳税，会影响人们持有什么。可是卡诺的贡赋大部分收布，市场上的日常货币却是贝币，可见纳税的规定并不单独决定人们用什么作钱。按巴特的说法，博尔努的贝币是出于统治者的投机而引进的；达荷美的国王则用限价与短串，直接从贝币上收税。

贝币最终让位给了国家的货币，但不是一纸禁令就做到的。Gardner 说，殖民前的欧洲商人愿意用贝币与布支付非洲的出口品，却不愿意收这些东西来卖枪支之类的其他进口品。原因是贝币与布在欧洲市场上价值很低，不能按同样的比价换成商行需要的欧洲货币。到 19 世纪末，殖民政府“开始立法禁止贝币与马尼拉铜环的进口，宣布它们不再是法偿货币”，再发行新的货币来取代它们\cite{gardner2015}。马尼拉铜环（manilla）是另一种当钱用的铜环。

禁令并没有立刻生效。Gardner 接着说，早期的殖民政府力量有限，这些法令往往难以执行，贝币等旧货币又用了几十年。让它们退出的还有别的力量。19 世纪的“贝币通胀”抬高了使用它们的成本，越来越多的非洲人领取英国货币的工资，与国际市场的往来也越来越多\cite{gardner2015}。贝币不必由国家发行；国家要废除它，也要等新货币在市场上用起来更方便。

\subsubsection{帝国能不能完全不用货币}
印加的证据支持“能”，但要打几个折扣。第一，它说的是国家的财政不用货币，不是没有交换。加西拉索自己说百姓以一种食物交换另一种食物，阿科斯塔也说美洲各地有大市场与著名的集市。第二，史料是征服之后写的。加西拉索晚年在西班牙写作，瓦莱拉的记载经他转引，两人都有为印加辩护的意思。“如此适度”的贡赋、按人口配置的记账人，都是理想化的图景。第三，这些记载描述的是制度应当怎样运作，没有给出它实际运作的数字。记载里没有一个印加仓库的库存数，仓储是否真的挡住了饥荒，也就无从检验。

Kocherlakota 的结论从理论上回答了这个问题。如果记忆足够完备，货币是多余的。印加的结绳官说明，这样的记忆可以用人力维持，但它只记录国家与每户之间的义务；百姓之间的大量往来，仍然要靠物物交换。

\begin{pitfall}[没有铸币，就没有货币，也没有市场]
这是两个误解叠在一起。第一，没有铸币不等于没有货币。贝币、盐、布、铜条与可可豆都承担了货币的职能，加奈姆的布还是记账单位。第二，没有货币不等于没有市场。阿科斯塔说，美洲“虽然有大市场与著名的集市，他们却不需要货币，也不需要经纪人，因为每个人都知道每样货物该换什么”\cite{acosta1880}。Gardner 批评过一种常见的看法，即“认为用布或贝壳的交换只是‘物物交换’，其中不用货币”\cite{gardner2015}。真正的物物交换的麻烦，出现在几种货币之间没有稳定比价的时候。巴特在库卡瓦买粮，要先把银元换成贝币，再把贝币换成衬衫，最后才换到粮食。
\end{pitfall}

\subsubsection{史料的偏差}
本章的数字都来自有立场的人。乌马里的资料是听来的。马克里齐的六迪拉姆在法译本里译错了，连法译者都不信，可见译本本身也会带来偏差。达尔泽尔与他的编者是奴隶贸易的辩护者，他们记下的达荷美国王的“暴政”，是从奴隶贩子的账本看出去的。

殖民时代与征服时代的观察者也各有立场。伯顿与罗宾逊是殖民时代的英国旅行者。罗宾逊是传教士，主张由英国出面废除奴隶贸易，认为第一步是找到一种令人满意的货币，来代替用作货币的奴隶；他说豪萨地区的问题在于“缺乏任何真正便携的货币”\cite{robinson1896}。加西拉索与瓦莱拉美化印加，阿科斯塔与奥维多是征服者一方的观察者。

数字之间的一致可以增加我们的信心。诺里斯看到的 1,280,000 枚贝币正好是 80 个贸易盎司，J. F. 算的重量与原注说的“约一吨半”相符；罗宾逊的脚夫工钱与他说的洛科贾价格一致。但一致只说明这些数字彼此相容，不说明它们代表了全貌。

\section{今天还剩什么}\label{sec:africa-legacy}

贝币、可可豆与结绳都早已不用了，但这一章的几个问题今天还在。

\textbf{货币是商品，还是国家的创造？}本章的证据两边都有。贝币、盐、布与可可豆都不是国家发行的，却都受到纳税规定与统治者的左右。门格尔（Carl Menger，1840—1921）把货币看成从交易中自发产生的商品，克纳普（Georg Friedrich Knapp，1842—1926）把它看成国家的创造，这场争论在\chref{ch:goldstandard}展开。现代货币理论沿着克纳普的一边走得更远（\chref{ch:minsky}），私人发行的稳定币则在另一边试探（\chref{ch:digitalmoney}）。

\textbf{货币与记账。}Kocherlakota 的“货币是记忆”把货币看成一种记账技术。今天的电子支付、央行数字货币与分布式账本，争论的正是由谁来记账、记下什么、谁能看到这本账（\chref{ch:digitalmoney}）。

\textbf{持有货币的成本。}可可豆的持有收益是负的；现金的收益是零减通胀。利率能降到零以下多少，取决于持有现金、把它存起来要花多少钱（\chref{ch:zlb}）。

\textbf{外来的货币与货币的等级。}西非的贝币供给由远方的产地与商人决定，殖民政府最终用西非镑这样的新货币取代了它，部分是为了让以英镑计的开支与债务不再承受汇率的风险\cite{gardner2015}。依赖别人决定的货币，今天表现为货币的等级与外围国家的金融依附（第~\ref{ch:periphery}、\ref{ch:hierarchy}~章）。

\textbf{制度的长期影响。}米塔制结束将近两百年后，它的边界仍然划分着今天秘鲁的消费与儿童健康。制度怎样影响长期发展，是\chref{ch:development}的问题。

\textbf{让财富留在本国。}阿达洪祖二世的理由，与同时代欧洲的重商主义者几乎一样，都是限制贸易，以便把财富留在国内。它在维达造成了什么，\chref{ch:mercantilism}还会在欧洲看到。

\begin{deeper}[单位、史料与几处不一致]
\textbf{单位。}米斯卡勒、迪拉姆、金第纳尔、杜卡特、英镑与磅，正文在第一次用到时已经说明，这里补充比较时的约定。本章把伊本·白图泰的第纳尔与巴特的米斯卡勒当作大致相同的单位，只用来比较数量级。迪拉姆与第纳尔的比价随时随地变化（\chref{ch:islamic}）。贝币：一串散的 40 枚、穿起来的 39 枚；J. F. 的一头为 4,000 枚，埃利斯的一头为 2,000 枚，Gardner 的“大头”为 4,000 枚。贸易盎司是记账单位，值金盎司的一半。阿科斯塔的四阿罗巴为一英担，estado 是一人高。

\textbf{“七百第纳尔的收益”。}乌马里的法译文 “pour trois cents dinars ils eurent sept cents dinars de gain” 可以读作利润七百，也可以读作连本带利七百；例~\ref{ex:africa-gold} 两种都算了，结论不变。

\textbf{伯顿的“头”。}伯顿在这段话里没有说一头是多少枚，本章按埃利斯的 2,000 枚折算。如果按 J. F. 的 4,000 枚，1863 年一先令换的贝壳要加倍，贬值就更快。伯顿自己算的一头值一先令九便士半，而 $54/2.5=21.6$ 便士，相差 0.1 便士，应是四舍五入。同一条注释接着给出一串与一美分各合多少，这两个折算彼此对不上，与前面的比价也对不上，无论一头按多少枚算都是如此，本章不用。

\textbf{两种重量。}J. F. 的贝壳每枚约 1.19 克，是原注所说的“实验”称出来的；罗宾逊的约 2.06 克，是“约一千磅”的估计。两者的差别与 Gardner 说的桑给巴尔贝壳更大相符，但罗宾逊的数字很粗。

\textbf{两个 1720 年以后的数字。}Gardner 说 1720 年 4,000 枚值约 1 英镑，J. F. 在 1793 年把 4,000 枚记作 10 先令。后者是按贸易盎司记的账，不一定是市场价格；两者之间的贬值可能有一部分来自记账的惯例。图~\ref{fig:cowrie_value}(b) 的各点来自不同的地方与不同的记账习惯，只能看出大致的趋势。

\textbf{可可豆能不能存。}阿科斯塔说它“可以久存而不坏”，马蒂尔说它“不能长久囤积”。两人都没有给出数字，例~\ref{ex:africa-cacao} 的损耗率是教学假设。

\textbf{巴特的两个年份。}巴特在卡诺说海岸上的贝币需求“1848 年以来”大增，在博尔努的脚注里说的是“1849 年以后”；两处说的是同一件事，本章两个年份都照录。

\textbf{马克里齐的六迪拉姆。}法译本两处都作第纳尔“跌到六迪拉姆”，法译者认为不可信，与乌马里相反。附录正文所译的那一处，阿拉伯原文说的是跌了六个迪拉姆\cite{vs2016}；按这个读法，马克里齐的跌幅至多 24\%，乌马里的至少 12\%，两人的记载可以相容。另一处（注释所引的马克里齐编年史）本章没有核对阿拉伯原文。本章采用乌马里的数字，因为他给出了前后两个价格。
\end{deeper}

\begin{takeaways}
\begin{itemize}
\item 14 世纪的西非，黄金、盐、铜、布与贝壳都承担货币的职能；黄金、盐与铜的开采与运输都离不开奴隶。价差受运费约束，式~\eqref{eq:africa-carry}：盐在马里城卖到伊瓦拉坦的 2—3.75 倍，黄金则可以一直运到开罗与欧洲。
\item 曼萨·穆萨 1324 年带来的黄金使开罗的金价至少下跌 12\%，十二年不回升。他向开罗商人借款，利率为 133\%—233\%，按式~\eqref{eq:africa-risk}，商人估计的还款概率只有三到五成；事后乌马里的消息来源说还了，伊本·赫勒敦说库维克家族的五万第纳尔没有收回。
\item 贝币不是国家发行的，它的价值由运费与供给决定：一枚贝壳在马里值马尔代夫的 348—1,043 倍；19 世纪桑给巴尔的贝壳大量涌入，一先令换的贝壳从 1720 年的 200 枚增加到 1894 年的 2,000 枚，平均每年 1.33\%。贝币运不动，大额支付就改用银元、布与奴隶。
\item 统治者可以左右贝币：达荷美的阿达洪祖二世以市价约一半的限价垄断奴隶贸易，又用每串短 3—6 枚的贝币付款，结果商路改道、商人破产、港口无贸易；博尔努的当权者用贝币投机，银元的价格高低相差 2.22 倍。
\item 印加帝国以劳役税与仓储代替货币，结绳是国家的“记忆”；仓储能提供多少保险取决于损耗率，式~\eqref{eq:africa-store}。西班牙人把劳役改造成采矿的米塔制（1573—1812），它的后果今天还在。可可豆是会坏的货币，持有它的收益为负，式~\eqref{eq:africa-hold}。
\end{itemize}
\end{takeaways}

\begin{selfcheck}
\begin{enumerate}
\item 一种布在甲地每匹卖 9 个单位，运到乙地卖 12 个单位，这条商路上一直有人贩运。运输的全部成本大约是甲地价格的多少？如果修了一条路，运费减半，乙地的布价会怎样变化？
\item 一位君主借款，债主要求 50\% 的利息，而别处的收益为 10\%。按式~\eqref{eq:africa-risk}，债主估计的还款概率是多少？如果君主把一座城的税收抵押给债主，利息降到 20\%，这个概率变成多少？抵押为什么能降低利率？
\item 设仓库每年结余为一年消费的 3\%。损耗率为 10\% 时，稳态库存是一年消费的百分之几？如果改进储存，损耗率降到 2\% 呢？
\item 设可可豆每年损耗 20\%，以可可豆计的物价下跌 5\%。持有可可豆一年的实际收益率是多少？
\item 一串贝币应穿 39 枚。如果国王付款时每串只穿 35 枚，相当于多少铸币税率？如果按散的 40 枚算呢？这与中世纪的货币贬值有什么相同之处？
\item 罗宾逊说，在扎里亚，好的坐骑卖十五先令到六英镑。按最高价 6 英镑、一先令 1,500 枚计，这匹马值多少枚贝币？按一磅 220 枚，重多少磅？要几个背 95 磅的脚夫？为什么在这样的市场上，奴隶与布会被用来缴纳大额的贡赋？
\end{enumerate}
\end{selfcheck}

\begin{answers}
\item 价差为 $12/9-1\approx33.3\%$；既然一直有人贩运，全部成本大致就是这么多。运费减半，乙地的布价会下跌（甲地可能略涨），直到价差降到新的成本附近。这就是式~\eqref{eq:africa-carry}。
\item $p=1.1/1.5\approx0.733$。利息降到 20\% 时，$p=1.1/1.2\approx0.917$。抵押让债主在君主违约时仍能收回一部分，相当于提高了收回的概率，债主要求的风险补偿就少了；曼萨·穆萨的债主没有这样的抵押。注意式~\eqref{eq:africa-risk} 假定违约时什么也收不回；有了抵押，算出的 0.917 应读作债主预期能收回的比例（连抵押在内），不一定是君主按时还款的概率。
\item $S^{*}=3/0.1=30$，即一年消费的 30\%；损耗率降到 2\% 时，$S^{*}=3/0.02=150$，即一年半。
\item $1+r^{m}=0.8/0.95$，$r^{m}\approx-15.8\%$。物价下跌使可可豆更值钱，但抵不过 20\% 的损耗。
\item $1-35/39\approx10.3\%$；按散的 40 枚算，$1-35/40=12.5\%$。两者都是付款的一方暗中减少每个单位所含的东西：中世纪的君主减少钱币的含银，达荷美的国王减少每串的枚数，都是式~\eqref{eq:medieval-seigniorage} 的铸币税。
\item $6\times20\times1\,500=180\,000$ 枚，约 $180\,000/220\approx818$ 磅，要 $818/95\approx8.6$ 个脚夫，即九个；最便宜的 15 先令也要 22,500 枚。用贝币缴纳大额贡赋要组织庞大的搬运队，而奴隶自己会走路，布每单位重量的价值也比贝币高；这正是罗宾逊说的豪萨诸邦以奴隶与布缴纳贡赋的原因。
\end{answers}

曼萨·穆萨带着一百驮黄金来到开罗，开罗的金价十二年没有回升。可是朝觐归来、回到开罗时，他却要向那里的商人借钱。按商人要的利息，他们把这位国王看成一个很可能赖账的借款人。库维克家族追到撒哈拉以南三次，一文也没有收回；腓力二世的银行家在停付之后，还能收回六到八成。一位君主怎样才能让债主相信他会还钱，从而以低利率借到大量的钱？下一章去看 17 世纪的荷兰与英格兰。
